\documentclass[a4paper,12pt]{article}
\usepackage[spanish,activeacute]{babel}
\usepackage[latin1]{inputenc}

%% Para las columnas
\usepackage{multicol}
\usepackage{amssymb,amsmath,eepic,fancyhdr}
\usepackage{latexsym}
\usepackage{datenumber}


%% Para numeraci?n autom?tica
\newcounter{nroproblema}
\setcounter{nroproblema}{1}
\newcounter{nroapartado}
\setcounter{nroapartado}{1}

\newcommand{\n}{%
    \vspace{.1in}\noindent{\bf \arabic{nroproblema}.\ }%
    \addtocounter{nroproblema}{1}%
    \setcounter{nroapartado}{1}%
    }%

\newcommand{\nn}{%
    \vspace{.1in}\noindent \hskip.4cm {\bf \arabic{nroproblema}.\ }%
    \addtocounter{nroproblema}{1}%
    \setcounter{nroapartado}{1}%
    }%


\newcommand{\apartado}{\vspace*{0.2cm}\hspace*{0.2cm}%
   {\bf \alph{nroapartado})}
   \addtocounter{nroapartado}{1}%
   }

% \newcommand{\A}{{\bf (A)\; }} % De problemas de Andrei
\newcommand{\A}{} % De problemas de Andrei

\newcommand{\R}{{\mathbb{R}}}
\newcommand{\C}{{\mathbb{C}}}
\newcommand{\N}{{\mathbb{N}}}
\newcommand{\K}{{\mathbb{K}}}


\newcommand{\rpol}[1]{\mathcal P^{\R}_{#1}}
\newcommand{\kpol}[1]{\mathcal P^{\K}_{#1}}

\newcommand{\cuatrox}[1][x]{\mbox{$\Bigg(\begin{array}{c}
    #1_1\\[-4pt] #1_2\\[-4pt] #1_3\\[-4pt] #1_4\end{array}\Bigg)$}}
\newcommand{\cuatroy}{\cuatrox[y]}
\newcommand{\cuatroz}{\cuatrox[z]}

% \newcommand{\cpol}[1]{\mbox{$\mathcal P^{\C}_{#1}}$}




%% m�rgenes

% Anterior:

%\setlength{\unitlength}{1mm} \addtolength{\voffset}{-15mm}
%\addtolength{\textheight}{30mm} \addtolength{\hoffset}{-20mm}
%\addtolength{\textwidth}{45mm}

\setlength{\unitlength}{1mm} \addtolength{\voffset}{-22mm}
\addtolength{\textheight}{40mm} \addtolength{\hoffset}{-22mm}
\addtolength{\textwidth}{46mm}




\begin{document}



% \begin{center}{\bf {\Large \'ALGEBRA I. HOJAS 2 -- 3 VERS 2 (PRELIMINAR)   }}\end{center}



\begin{flushright}
1${}^o$ de grado en F�sica, curso 2013/14

\today
\end{flushright}


\begin{center}{\bf {\Large \'ALGEBRA I. HOJA 4}}\end{center}

 \vskip.35cm


%\noindent{Nota: Las hojas est\'an colgatas tambi\'en en la p\'agina web de del curso
%\begin{verbatim}http://www.uam.es/personal_pdi/ciencias/mzambon/WS12.html\end{verbatim}}



% \begin{center}{\bf {\Large \'ALGEBRA I. PARA LAS HOJAS POSTERIORES}}\end{center}



%\n Estudiar si son subespacios vectoriales de $\mathbb{R}^2$ los siguientes subconjuntos:
%
%\apartado $S_1=\{(x, y) \in \mathbb{R}^2,\, y -x = 0\}$,
%
%\apartado $S_2 =\{(x, y) \in \mathbb{R}^2,\, y- x = 1\}$,
%
%\apartado $S_3 =\{(x, y) \in \mathbb{R}^2,\, xy = 0\}$,
%
%\apartado $S_4 =\{(x, y) \in \mathbb{R}^2,\, y = |x|\}$,
%
%\apartado $S_5 =\{(x, y) \in \mathbb{R}^2,\, |y| = |x|\}$,
%
%\apartado $S_6=\{(x,y) \in \mathbb{R}^2,\, x\geq 0 \}$.
%
%\vspace{0.2cm}



\n Sea $V$ un espacio vectorial.

\apartado Suponiendo que $v_1, v_2$ es una familia de vectores
linealmente independiente, demostrar que la familia
$2v_1+v_2, 3v_1+2v_2$ tambi�n lo es.

\apartado Supongamos ahora que $v_1, v_2, v_3$ es un sistema de generadores de $V$.
Demostrar que \newline $v_1+v_2, v_2+v_3, v_1+v_3$ forman tambi�n un sistema de generadores de $V$.

\n Sea $V$ un espacio vectorial y $A$ una transformaci�n lineal de $V$ a $V$.
Dado un polinomio $p(t)=\sum_{k=0}^n a_k t^{n-k}$, ponemos $p(A)=\sum_{k=0}^n a_k A^{n-k}$
(donde $A^0:=I_V$). Suponiendo que $p(0)=a_n=0$ y que $p(A)=I_V$, demostrar que
$A$ es invertible.

{\it Sugerencia:} Considerar primero el caso de un polinomio concreto, por
ejemplo, escoger $p(t)=t^4-t$.



%\n Sea $V$ un espacio vectorial y $A$ una transformaci�n lineal de $V$ a $V$.
%Supongamos que $A^n=I_V$ para alg�n n�mero natural $n>1$. Demostrar que $A$ es invertible.



\n Estudiar si son subespacios vectoriales del espacio vectorial real de funciones reales
de variable real los siguientes subconjuntos:

\apartado $S_1=\{f:\mathbb{R}\to\mathbb{R}  : \; f (0) = 0\}$; \qquad \qquad
\apartado $S_2 = \{f:\mathbb{R}\to\mathbb{R} : \; f (1) = 0\}$;
\vskip-.2cm

\apartado $S_3 = \{f:\mathbb{R}\to\mathbb{R} : \; f (0) = 1\}$.


%\n Averiguar si son subespacios vectoriales de los correspondientes espacios vectoriales los
%siguientes subconjuntos:
%
%\apartado  El subconjunto de los polinomios de variable real y con coeficientes reales, dentro de las
%funciones reales de variable real.
%
%\apartado El subconjunto de los polinomios de grado $\leq n$, de variable real y con coeficientes reales,
%siendo $n$ fijo, dentro del espacio vectorial de las funciones reales de variable real.

%\apartado El subconjunto de polinomios divisibles por $x-1$,  grado $\leq 3$,  de variable real y con
%coeficientes reales,  como subconjunto del espacio vectorial  de los polinomios de variable real,
% grado $\leq 3$, y con coeficientes reales.



\n Averiguar si son subespacios vectoriales de $\mathcal{M}_{2\times 2}(\mathbb{R})$ los siguientes subconjuntos:

\apartado Las matrices $2\times 2$ con n\'umeros racionales.

\apartado Las matrices de n\'umeros reales de orden $2\times 2$ de traza cero. (Se llama traza de una
matriz a la suma de los elementos de su diagonal).

\apartado Las matrices de n\'umeros reales de orden $2\times 2$ de rango menor o igual que $1$
(son las matrices, cuyas dos columnas son linealmente dependientes).

\apartado Las matrices de n\'umeros reales de orden $2\times 2$ que conmutan con la matriz
B, siendo B una matriz fija $2\times 2$.

% \vspace{0.2cm}

\n Averiguar si los siguientes subconjuntos son subespacios vectoriales de $\mathcal{M}_{n\times n}(K)$:

\apartado Las matrices diagonales.

\apartado Las matrices triangulares superiores.

\apartado Las matrices sim\'etricas (las que satisfacen que $A^T=A$).

\apartado Las matrices antisim\'etricas (las que satisfacen que $A=-A^T$).

% \vspace{0.2cm}

\vskip-.2cm






\n
{\bf a)}
Halla
$T\small\begin{pmatrix} 1 \\ 0\end{pmatrix}$ si $T:\R^2\longrightarrow\R^2$ es una aplicaci{\'o}n
lineal de la que sabemos que
\vskip-.3cm
\[
T
\small
\begin{pmatrix} 3 \\ 1
\end{pmatrix}
=
\begin{pmatrix} 1 \\ 2
\end{pmatrix}, \qquad
T
\small
\begin{pmatrix} -1 \\ 0
\end{pmatrix}
=
\begin{pmatrix} 0 \\ 1
\end{pmatrix}\, .
\]
%
%$T(3,1)=(1,2)$ y $T(-1,0)=(1,1)$.
%

{\bf b)} Lo mismo sabiendo que \enspace
$T
\small
\begin{pmatrix} 4 \\ 1
\end{pmatrix}
=
\begin{pmatrix} 1 \\ 1
\end{pmatrix}, \;
T
\small
\begin{pmatrix} 1 \\ 1
\end{pmatrix}
=
\begin{pmatrix} 3 \\ -2
\end{pmatrix}\, .
$
%
%$T(4,1)=(1,1)$ y $T(1,1)=(3,-2)$.
%


% {\bf PRODUCTOS ESCALARES Y NORMAS MINIMAS:}


\n Sea $M_{m,n}(\mathbb{R})$ el espacio de matrices $m\times n$ sobre $\mathbb{R}$,
sea $B^T$ la matriz traspuesta de $B$, y sea  $tr (A)$ la
traza de la matriz cuadrada  $A = (a_{ij})_{i,j = 1}^n$, definida mediante
$tr (A) := \sum_{i=1}^n a_{ii}$.
Probar
que la traza define un producto escalar en $M_{m,n}(\mathbb{R})$ mediante $(A,B):= tr(B^T A)$.

\n Sean $A,B$ matrices de tama�os $n\times m$ y $m\times n$, respectivamente.
Demostrar que $tr(AB)=tr(BA)$.

\n Probar que $\left\{\left( \begin{matrix}1 & 0 \\
0 & 0\end{matrix} \right), \left( \begin{matrix}0 & 1 \\
0 & 0\end{matrix} \right), \left( \begin{matrix}0 & 0 \\
1 & 0\end{matrix} \right), \left( \begin{matrix}0 & 0 \\
0 & 1\end{matrix} \right)\right\}$ es una base ortonormal de
$M_{2,2}(\mathbb{R})$, con producto escalar $(A,B)= tr(B^T A)$.





\n Sea $P(\mathbb{R})$ el espacio de polinomios  sobre $\mathbb{R}$. Probar
que $(f,g):= \int_0^1 f(t) g(t) dt$ define un producto escalar en $P(\mathbb{R})$.



\n
\apartado Comprobar que $\big((a,b)^T,(c,d)^T\big)_1:= ac - 2 ad - 2bc + 5bd$ define un producto escalar
en el plano real.

\apartado Calcular las normas de los vectores $u=(1,2)^T$  y $v=(4,5)^T$
con respecto al producto escalar habitual $(\cdot, \cdot)_h$  y con respecto al producto
$(\cdot, \cdot)_1$. Calcular tambi�n $(u, v)_h$ y $(u, v)_1$.


\n Hallar el �ngulo entre los vectores $e_1$ y $(1,1,1)^T$ en $\mathbb{R}^3$.

\n \apartado Sean $u= (z_1, z_2)^T$ y $v = (w_1, w_2)^T$ vectores en $\mathbb{C}^2$. Comprobar que
\[
(u,v)_1:=
z_1 \overline{w}_1  + (1+i)z_1 \overline{w}_2 + (1-i)z_2 \overline{w}_1 + 3 z_2 \overline{w}_2
\]
define un producto herm�tico en~$\mathbb{C}^2$.

\apartado Calcular las normas de los vectores $u=(1 - i,2 + 3i)^T$  y $v=(4 + 2i, 5 + i)^T$
con respecto al producto escalar habitual $(\cdot, \cdot)_h$  y con respecto al producto $(u,v)_1$.
Calcular tambi�n $(u, v)_h$ y $(u, v)_1$.






\n Decidir razonadamente si la siguiente afirmaci�n es verdadera o falsa:
$$
\bigg(\int_{99 \pi/200}^{\pi/2} \sen x\cos x dx\bigg)^2 \le \int_{99 \pi/200}^{\pi/2} \left(\sen x\right)^2 dx
\cdot \int_{99 \pi/200}^{\pi/2} \left(\cos x \right)^2dx.
$$
\vskip-.2cm

\n  Sea L la recta que pasa por
$(-2, 1, 3)^T$ y   $(1, 2, 4)^T$. Hallar el punto de L m�s pr�ximo al origen, y la distancia
m�nima. \enspace
\textit{Respuestas:} $(-16/11, 13/11, 35/11)^T, \sqrt{150/11}.$


\n
Escribiendo $y = y(x)$ (de modo que $x$ es la variable
independiente), y con los datos siguientes, ha\-llar la
correspondiente recta de regresi�n $y = a x + b$.

Datos: $(x,y) = (1,1), (3,2), (4,4), (6,4), (8,5), (9,7), (11, 8),
(14, 9)$.


\n
Encontrar el polinomio $y = ax^2+bx+c$
que mejor se ajusta, en el sentido de los m�nimos cuadrados,
a los siguientes datos:
 $(x,y) = (-2,0), (-1,0), (0,0), (1,1), (2,4)$.
Comparar gr�ficamente estos datos con
el polinomio.


% Soluci�n: $y(x)=(9x+5x^2)/10$.






\n Escribir la matriz de la simetr\'ia (reflexi�n) de $\mathbb{R}^3$ respecto a:

\apartado El plano de ecuaci\'on y = x. \qquad \qquad \qquad
\apartado El plano de ecuaci\'on y = z .
\vskip-.3cm

\apartado El plano de ecuaci\'on x = z . \qquad \qquad \quad\;\,
\apartado La recta de ecuaciones y = x, z = 0.
\vskip-.3cm

\apartado La recta de ecuaciones y = z , x = 0. \qquad
\apartado La recta de ecuaciones x = z , y = 0.




\n Hallar una base para el complemento ortogonal del subespacio generado por $(1,2,3)$ en $\mathbb{R}^3$.


\n
\apartado Sea $P_2(\mathbb{R})$ el espacio de polinomios sobre $\mathbb{R}$ de grado $\le 2$, con
el  producto escalar
$(f,g)_1:= \int_0^1 f(t) g(t) dt$. Hallar una base para el complemento ortogonal
del subespacio generado por $1 + 3t$.

\apartado Hacer lo mismo para el caso cuando el producto escalar en $P_2(\mathbb{R})$ se define por
$(f,g)_2:= \int_{-1}^1 f(t) g(t) dt$.
\vskip.2cm

?`Coinciden los complementos ortogonales en estos
dos casos?




\end{document}
